\chapter*{Tóm tắt}
\label{abs}

Học sâu hiện đại phụ thuộc nặng vào dữ liệu lớn, khiến các bài toán thiếu dữ liệu gán nhãn gặp nhiều trở ngại. Trong bối cảnh đó, meta-learning dựa trên gradient, tiêu biểu là MAML, nổi lên như hướng tiếp cận chủ đạo cho học ít mẫu. Tuy nhiên, cơ chế tri thức tiền nghiệm ở tham số khởi tạo $\theta_0$ chi phối quá trình thích ứng nhanh vẫn là hộp đen, khiến chuyển giao tiêu cực (negative transfer) - khi tập hỗ trợ $S_i$ định hướng sai làm giảm hiệu suất — chỉ được phát hiện qua các biểu hiện chứ chưa giải thích được ở cấp độ nguyên nhân. Đề tài đặt mục tiêu xây dựng phương pháp hậu diễn giải nhằm định lượng mức lợi/hại của bước thích nghi và định vị đặc trưng trong tập hỗ trợ gây ra ảnh hưởng đó. Trong đó, đề tài đề xuất độ lợi thích nghi $\Delta M$, đo bằng chênh lệch hàm mất mát trên tập truy vấn giữa hai nhánh can thiệp và không can thiệp quá trình thích ứng, cùng phương pháp \acrfull{fama} nhằm mở rộng Class Activation Mapping qua nhiều bước gradient để trực quan hóa đặc trưng có lợi/hại. Tính trung thực được kiểm định qua phép xóa-và-đo hai chiều và kiểm tra tỉnh táo, thực nghiệm trên kiến trúc Conv4 với \emph{miniImageNet}, \emph{tieredImageNet} và \emph{CUB-200-2011}. Kết quả cho thấy hai chỉ số $\mathrm{ABC}_\downarrow$, $\mathrm{ABC}_\uparrow$ đạt ý nghĩa thống kê (p<0,001), tương quan tỉnh táo sụt giảm gần như tức thời khi phá hủy tầng gần đầu ra, và phân tích trên 600 tác vụ chỉ ra "thiên lệch ngữ cảnh" là nguyên nhân cốt lõi gây chuyển giao tiêu cực. Kết quả khẳng định \acrshort{fama} là công cụ diễn giải trung thực, đáng tin cậy cho thích nghi trong meta-learning, mở ra ứng dụng chẩn đoán, giảm thiểu chuyển giao tiêu cực, và làm nền tảng mở rộng sang kiến trúc sâu hơn.
